\chapter{赛题需求分析}
\label{chap:introduction}

\section{研究背景与应用价值}
本项目面向 2026 年全国大学生嵌入式芯片与系统设计竞赛 FPGA 创新设计赛道，
选择紫光同创选题二“基于紫光同创 FPGA 的任意角度的畸变矫正”。
该赛题要求利用国产 FPGA 平台构建完整的视频处理系统，在真实相机标定参数的基础上，
实现图像几何矫正，并进一步结合姿态感知、硬件流水线及自适应增强，
提升动态场景下的实时处理能力与成像效果\cite{problem}。

在机器视觉、自动驾驶、智能交通和视觉测量等应用中，图像不仅用于观察，
也是目标定位与空间测量的数据来源。镜头的径向畸变可能使直线呈现桶形或枕形弯曲，
切向畸变则会造成局部位置偏移；相机倾斜和转动还会改变目标的投影关系。
这些因素会影响图像拼接、目标检测及尺寸测量的准确性。
因此，需要将镜头畸变补偿与相机姿态变换结合起来，使输出图像保持符合预期的几何关系。

高分辨率视频的几何矫正需要逐像素进行坐标变换与重采样，
且变换后的源像素访问通常不连续。FPGA 可将坐标计算、图像缓存和视频接口组织为
并行模块与流水线，为持续视频处理提供可控的时序基础。
其实现难点在于有限逻辑、乘法和片内存储资源下，协调运算精度、外部存储带宽与处理延迟。
本项目以可改变姿态的相机为应用对象，通过云台调整拍摄方向、IMU 提供姿态信息，
在 FPGA 上完成图像矫正与增强，并由上位机观察输出和控制功能，
形成可用于动态演示与客观测评的实验系统。

\section{赛题任务分解}
依据选题指南正文第 10--12 页，本系统的任务分为基础功能与高阶功能两部分。
基础功能构成从视频输入到矫正输出的数据链路，高阶功能进一步解决动态姿态、
实时吞吐和复杂光照条件下的适应性问题\cite{problem}。

\subsection{基础功能需求}
\begin{enumerate}[label=（\arabic*）]
  \item \textbf{图像采集、缓存与显示。}
  实现 HDMI 或摄像头接口的视频输入、图像缓存和显示输出，保证像素数据与同步信号一致。
  本系统采用 HDMI 输入输出和 DDR3 帧缓存，为几何变换提供可寻址的源图像。
  该链路还需协调输入、存储和输出时钟域，避免读写冲突及画面撕裂。

  \item \textbf{畸变模型与坐标映射。}
  在 FPGA 中实现二阶径向及切向畸变模型，依据焦距、主点和畸变系数计算像素对应关系。
  标定参数应来源于真实相机，并保持标定分辨率、像素坐标约定与硬件实现一致。
  本系统采用输出像素驱动的逆映射，由理想输出位置反求畸变源图像中的采样坐标。

  \item \textbf{图像重采样与完整性。}
  根据映射得到的非整数坐标进行重采样，赛题给出双线性插值作为基础算法示例，
  并要求消除矫正后的黑边或空洞现象。
  设计中需分别处理离散映射产生的漏采样与变换后超出源视场的区域：
  逆映射可避免前者，后者需通过有效视场选择、裁切等策略控制，并评价其对视场的影响。
\end{enumerate}

\subsection{高阶功能需求}
\begin{enumerate}[label=（\arabic*）]
  \item \textbf{任意角度与动态矫正。}
  结合 IMU 或上位机指令，在相机倾斜、旋转时动态更新几何变换。
  姿态坐标轴、角度符号和安装零点应与图像坐标系保持一致；
  同一帧处理期间应采用一致的姿态参数，避免参数切换引起局部错位。
  动态验证需覆盖连续角度变化及不同测试图卡，记录支持的角度范围与边界行为。

  \item \textbf{硬件加速与流水线优化。}
  针对坐标映射中的乘除法、定点运算和非连续访存设计硬件流水线，
  通过缓存复用和突发传输提高外部存储效率。
  系统吞吐应包含采集、运算、访存和结果发布全过程，
  并满足持续视频处理所需的帧率与时序约束。

  \item \textbf{自适应图像预处理。}
  面向强光、夜间等复杂光照，在矫正前或矫正过程中加入自适应 Gamma 校正
  或直方图均衡等增强算法。本系统选择帧自适应 Gamma 方案，
  根据亮度统计调整映射曲线，并通过增强开关对比暗部细节、高光表现及帧间稳定性。
\end{enumerate}

\section{设计目标与评价指标}
\subsection{系统规格与实时处理目标}
赛题以 1080P@30fps 或 60fps 为实时处理标准的示例，
同时要求测量端到端硬件延迟，并在满足性能的前提下降低资源占用。
结合所选 PGL50H 器件与帧缓存结构，本系统将输入输出配置为
$1920\times1080$、60 Hz 扫描，将完整矫正结果的更新速率设为不低于 30 fps 的工程目标。
按有效图像区域计算，30 fps 对应每秒处理 62.208 百万个输出像素，
每个结果帧的平均时间预算约为 33.33 ms。

显示扫描率与矫正更新率分别统计：前者反映输出接口的扫描时序，
后者统计单位时间内由不同源帧生成的完整结果帧数量。
重复显示上一结果帧不计入新的矫正帧，单帧渲染耗时也不等于端到端延迟。
端到端硬件延迟以 FPGA 视频输入侧接收某一可识别事件为起点、
视频输出侧输出其对应像素为终点；若使用相机和采集卡测量整条演示链路，
还需说明这些设备引入的额外延迟。

\subsection{算法精度与工程评价方法}
表~\ref{tab:requirements} 给出本系统依据赛题建立的评价项目。
赛题对清晰度、直线度和延迟给出了评价方向，但未规定统一的数值门限，
因此本报告采用固定图卡、固定测试区域和可重复的测量方法记录结果，
并结合软件参考与原始图像进行比较。

\begin{table}[htbp]\centering\small
\caption{系统设计目标与评价方法}\label{tab:requirements}
\begin{tabularx}{\textwidth}{>{\raggedright\arraybackslash}p{2.5cm}YY}
\toprule
评价项目 & 设计目标或评价重点 & 测量与判定方法 \\
\midrule
视频链路 & $1920\times1080$，60 Hz 输入输出扫描，RGB888 接口 & 检查有效像素数、同步时序、颜色和帧完整性 \\
标定与坐标精度 & 正确使用真实相机内参及畸变系数，控制定点误差 & 核对参数来源，与独立浮点参考比较最大及均方根坐标误差 \\
图像质量 & 直线无明显弯曲，边缘无明显拉伸或断裂 & 使用网格和斜边图卡，测量直线度及 MTF50 或 SFR，记录有效视场与无效区域比例 \\
实时性与延迟 & 完整矫正结果更新率 $\geq30$ fps，并降低链路延迟 & 按源帧身份统计发布率，分别测量渲染耗时和端到端延迟 \\
动态姿态 & 在声明的角度范围内连续更新，帧内参数一致 & 测试中间及边界角度、连续转动、零点和异常姿态输入 \\
自适应增强 & 改善复杂光照下的可见细节，控制亮度跳变 & 对比增强开关及亮暗切换场景，记录亮度统计和主客观效果 \\
资源与稳定性 & 资源占用不超过器件容量，满足时序并保持连续运行 & 检查 LUT、FF、APM、DRM 报告、CDC 处理及持续运行中的异常计数 \\
\bottomrule
\end{tabularx}
\end{table}

\section{开发约束与验证要求}
硬件平台采用赛题推荐的 MES50HP 开发板，核心处理器件为紫光同创 PGL50H，
使用 Verilog HDL 开发，采用 Pango Design Suite（PDS）完成综合、布局布线、
时序分析和位流生成。标定与测试环境由参赛者自行准备，
需包含真实相机、棋盘格或圆点标定板、标定软件及上位机图像质量分析工具。
赛题列举 OpenCV 标定工具箱用于生成基准参数，本工程使用实际相机的 MATLAB 标定记录，
后续比较应统一参数含义、坐标约定及畸变模型。

工程实现需重点处理视频输入、存储控制、姿态接收和显示输出之间的跨时钟域通信，
合理分配片内存储并约束外部接口时序。功能划分、接口定义、关键定点格式和时序约束
应在代码及文档中保持一致。验证采用算法参考、RTL 仿真、FPGA 实现分析与板级测试
相结合的方式，保存与源码、约束和位流版本对应的记录。
决赛演示应能够连续输入实时视频，并在改变相机物理角度或更换图卡时展示输出效果。

\section{本报告的工作内容}
本报告围绕“视频采集与缓存、姿态与镜头模型、逐像素逆映射、图像重采样及显示”
组织系统设计，重点说明定点运算流水线、分块访存和帧级参数管理之间的配合。
第~\ref{chap:system}~章介绍开发平台、硬件组成与总体数据通路；
第~\ref{chap:algorithm}~章介绍标定参数、姿态变换、畸变模型及重采样方法；
第~\ref{chap:hardware}~章说明 FPGA 模块划分、存储调度与跨时钟域设计；
第~\ref{chap:control}~章介绍姿态通信、上位机控制及自适应 Gamma；
第~\ref{chap:verification}~章给出验证方案和板级测试方法；
第~\ref{chap:performance}~章分析帧率、延迟、资源及图像质量，最后总结局限与改进方向。

当前工程基线采用最近邻重采样，yaw/pitch 每轴支持范围约为 $\pm15^\circ$，
roll 未参与变换，所用标定的切向系数为零，无效源区域输出黑色。
因此，双线性插值的实时实现、非零切向场景、完整姿态自由度及黑边处理仍需单独评估。
已有归档提供 RTL 仿真和布局布线证据，实板画质与端到端延迟尚需测量；
后续各章据此说明实现范围，并以相应测试记录判定赛题要求的满足情况\cite{project,current}。
